\documentclass[11pt,a4paper]{article}
\usepackage[utf8]{inputenc}
\usepackage[T1]{fontenc}
\usepackage[brazil]{babel}
\usepackage{lmodern}
\usepackage[margin=2cm]{geometry}
\usepackage{amsmath,amssymb}
\usepackage{booktabs}
\usepackage{graphicx}
\usepackage{hyperref}
\usepackage{xcolor}
\usepackage{titlesec}
\usepackage{tabularx}
\usepackage{enumitem}
\usepackage{float}
\usepackage{array}
\renewcommand{\tabularxcolumn}[1]{>{\raggedright\arraybackslash}p{#1}}
\usepackage{newunicodechar}
\usepackage[most]{tcolorbox}
\usepackage{etoolbox}

\hypersetup{colorlinks=true, linkcolor=blue!70!black, citecolor=blue!70!black, urlcolor=blue!70!black}

\titleformat{\section}{\large\bfseries}{\thesection.}{0.5em}{}
\titleformat{\subsection}{\normalsize\bfseries}{\thesubsection.}{0.5em}{}
\titlespacing*{\section}{0pt}{1.2ex plus 0.5ex minus .2ex}{0.6ex}
\titlespacing*{\subsection}{0pt}{1.0ex plus 0.4ex minus .2ex}{0.4ex}
\setlength{\parskip}{0.45em}
\setlength{\parindent}{0pt}
\setlist{itemsep=0.15em,topsep=0.3em}

% --- símbolos ---
\newunicodechar{π}{\ensuremath{\pi}}
\newunicodechar{τ}{\ensuremath{\tau}}
\newunicodechar{φ}{\ensuremath{\varphi}}
\newunicodechar{ψ}{\ensuremath{\psi}}
\newunicodechar{Σ}{\ensuremath{\Sigma}}
\newunicodechar{→}{\ensuremath{\rightarrow}}
\newunicodechar{⇒}{\ensuremath{\Rightarrow}}
\newunicodechar{∈}{\ensuremath{\in}}
\newunicodechar{−}{\ensuremath{-}}
\newunicodechar{≈}{\ensuremath{\approx}}
\newunicodechar{≠}{\ensuremath{\neq}}
\newunicodechar{≡}{\ensuremath{\equiv}}
\newunicodechar{≤}{\ensuremath{\leq}}
\newunicodechar{≥}{\ensuremath{\geq}}
\newunicodechar{′}{\ensuremath{'}}
\newunicodechar{‑}{-}
\newunicodechar{✔}{\ensuremath{\checkmark}}
\newunicodechar{✘}{\ensuremath{\times}}

% --- cores e caixas ---
\definecolor{cC}{HTML}{7A4FA3}
\definecolor{cS}{HTML}{2F7D74}
\definecolor{cAcc}{HTML}{2F5D8A}
\newcommand{\metrictag}[2]{{\setlength{\fboxsep}{1.5pt}%
  \ifstrequal{#1}{c}{\colorbox{cC}{\textcolor{white}{\scriptsize\textbf{#2}}}}{%
  \ifstrequal{#1}{s}{\colorbox{cS}{\textcolor{white}{\scriptsize\textbf{#2}}}}{%
  \colorbox{black!60}{\textcolor{white}{\scriptsize\textbf{#2}}}}}}}
\newcommand{\pagechip}[1]{{\setlength{\fboxsep}{1.5pt}\colorbox{cAcc}{\textcolor{white}{\scriptsize\texttt{#1}}}}}
\newcommand{\readtag}[1]{\texttt{\small #1}}
\newcommand{\partbar}[1]{\par\vspace{1.2em}\noindent\colorbox{black!80}{\parbox{\dimexpr\linewidth-2\fboxsep}{\color{white}\bfseries\sffamily\MakeUppercase{#1}}}\par\vspace{0.4em}}
\newtcolorbox{citacard}[1]{breakable, enhanced, colback=white, colframe=black!25, boxrule=0.5pt, arc=2pt,
  left=6pt, right=6pt, top=4pt, bottom=4pt, title={\small #1}, coltitle=black, colbacktitle=black!6, fonttitle=\normalfont}
\newenvironment{origquote}{\par\itshape\small}{\par}
\newenvironment{tradquote}{\par\small\textsc{Tradução:}\ }{\par}
\newcommand{\explab}[1]{\par\smallskip{\small\color{cAcc}\textbf{\MakeUppercase{#1}}}\par}
\colorlet{notebg}{blue!4}\colorlet{notefg}{cAcc}
\colorlet{okbg}{green!5}\colorlet{okfg}{green!45!black}
\colorlet{warnbg}{orange!8}\colorlet{warnfg}{orange!80!black}
\newtcolorbox{infobox}[2][note]{breakable, enhanced, boxrule=0pt, arc=2pt, left=6pt, right=6pt, top=4pt, bottom=4pt,
  colback=#1bg, borderline west={3pt}{0pt}{#1fg}, frame hidden,
  before upper={\ifstrempty{#2}{}{\textbf{#2}\par}}}
\colorlet{pairc}{cC}\colorlet{pairs}{cS}
\newenvironment{paircol}[2]{\begin{tcolorbox}[width=0.485\linewidth, nobeforeafter, box align=top, boxrule=0.5pt,
  colframe=black!25, colback=white, borderline north={3pt}{0pt}{pair#1}, left=5pt, right=5pt,
  before upper={{\color{pair#1}\small\textbf{\MakeUppercase{#2}}}\par}]}{\end{tcolorbox}}

\newunicodechar{ℓ}{\ensuremath{\ell}}
\newunicodechar{⊙}{\ensuremath{\odot}}
\newunicodechar{α}{\ensuremath{\alpha}}
\newunicodechar{β}{\ensuremath{\beta}}

\begin{document}

\begin{center}
    {\LARGE \textbf{Fichamento: disparidades na qualidade das explicações}} \\[0.35em]
    {\large DAI, J.; UPADHYAY, S.; AÏVODJI, U.; BACH, S. H.; LAKKARAJU, H. \textit{Fairness via Explanation Quality: Evaluating Disparities in the Quality of Post hoc Explanations}. AAAI/ACM Conference on AI, Ethics, and Society (AIES), 2022} \\[0.3em]
    \textbf{Tema:} diferenças de fidelidade, estabilidade, consistência e esparsidade das explicações entre grupos \\[0.2em]
    \textbf{Aluno:} Arthur Santos de Oliveira (RA: 156297) \\[0.2em]
    \textbf{Atualizado em:} 28/10/2026 \\[0.2em]
    \textbf{Link:} \url{https://doi.org/10.1145/3514094.3534159} (arXiv: \url{https://arxiv.org/abs/2205.07277}) \\[0.2em]
    \textbf{Leitura:} íntegra (arXiv v2, 12 páginas) \\[0.2em]
    \textbf{Repositório:} \url{https://github.com/Arthur-so/shap-ecg-arrhythmia}
\end{center}
\vspace{-0.5em}
\hrule
\vspace{0.8em}

\textbf{Como ler.} \texttt{[p. N]} indica a página do PDF lido (arXiv v2). As citações estão no original em inglês, cada uma seguida de uma tradução livre e de uma explicação.

\begin{infobox}[note]{Resumo}
\begin{itemize}
  \item \textbf{Objetivo:} propor uma estrutura para detectar se a qualidade das explicações difere entre grupos da população (sexo, raça).
  \item \textbf{Dados e modelos:} German Credit, Student Performance e COMPAS; regressão logística (LR) e rede neural de 3 camadas (NN).
  \item \textbf{Métodos de explicação:} LIME, SHAP, MAPLE, SmoothGrad, Integrated Gradients e gradiente simples.
  \item \textbf{Avaliação:} quatro propriedades medidas por instância (fidelidade, estabilidade, consistência, esparsidade), médias por grupo e teste de Mann-Whitney entre os grupos; 5 repetições com divisões diferentes dos dados.
  \item \textbf{Conclusão:} disparidades aparecem com frequência, mais em modelos complexos e mais com Integrated Gradients e SHAP; a disparidade em uma métrica não indica disparidade nas outras.
\end{itemize}
\end{infobox}

\section{O que o trabalho faz}
O artigo distingue duas questões: se a explicação retrata corretamente a justiça do modelo (tema de trabalhos sobre \textit{fairwashing}) e se o método de explicação funciona igualmente bem para todos os grupos. Ele trata da segunda. Escolhe quatro propriedades de qualidade, calcula cada uma por instância, separa as instâncias de teste pelo atributo sensível, compara as distribuições dos dois grupos com um teste de hipótese e discute o que cada tipo de disparidade significaria na prática (por exemplo, um médico confiando mais no modelo para um grupo do que para outro).

\section{Como avalia as explicações}
\begin{itemize}
  \item \textbf{Fidelidade:} (a) \textit{ground truth fidelity}, fração das $k = 5$ características mais importantes da explicação que coincidem com as do modelo linear; (b) \textit{prediction gap fidelity}, mudança média da probabilidade predita quando se soma ruído gaussiano ($\sigma = 0{,}1$, $m = 1.000$ sorteios) às características \textbf{fora} das $k$ principais (quanto menor, melhor).
  \item \textbf{Estabilidade:} distância L1 média entre a explicação de $x$ e as de vizinhos ruidosos ($m = 5$), como em Yeh et al.
  \item \textbf{Consistência:} distância L1 média entre explicações do mesmo ponto calculadas com sementes diferentes ($m = 5$).
  \item \textbf{Esparsidade:} número de características com importância acima de $t = 0{,}01$.
  \item \textbf{Disparidade:} teste U de Mann-Whitney entre as distribuições dos dois grupos.
\end{itemize}

\section{Resultados e conclusões}
\begin{center}
\small
\begin{tabularx}{\linewidth}{Xl}
\toprule
Achado & Valor relatado \\
\midrule
Combinações (método, métrica, modelo, dados) com disparidade significativa & 18,5\% de 162 \\
Combinações método--modelo--dados com disparidade em ao menos uma métrica & 56\% de 36 \\
Disparidade de fidelidade (\textit{prediction gap}) & 30,6\% (63,6\% delas na NN) \\
Disparidade de esparsidade / estabilidade / consistência & 30,6\% / 13,9\% / 1 de 36 \\
Por método: IntGrad, SHAP, gradiente, LIME, SmoothGrad, MAPLE & 33,3\%, 29,6\%, 22,2\%, 11,1\%, 11,1\%, 3,7\% \\
\bottomrule
\end{tabularx}
\end{center}
As disparidades de fidelidade e de estabilidade são mais frequentes e maiores no modelo não linear; nele também as variâncias são maiores, o que, segundo os autores, torna essas métricas menos robustas em modelos complexos. Integrated Gradients é o único método com disparidade em pelo menos uma métrica em todas as combinações de modelo e dados. Nas combinações com disparidade, metade a apresenta em uma só métrica e metade em duas.

\section{Citações}

\begin{citacard}{\pagechip{p. 5} §3.4.4 Disparidade de esparsidade}
\begin{origquote}“Unlike the three metrics discussed previously, disparity in sparsity across demographic groups is not inherently problematic. In fact, it may reflect the model’s true behavior, e.g., if the model has a more complex decision boundary for one group than another, or if the model does rely more features for one group than another.”\end{origquote}
\begin{tradquote}“Ao contrário das três métricas discutidas antes, a disparidade de esparsidade entre grupos demográficos não é problemática em si. Na verdade, ela pode refletir o comportamento real do modelo, por exemplo, se o modelo tem uma fronteira de decisão mais complexa para um grupo do que para outro, ou se de fato depende de mais características para um grupo do que para outro.”\end{tradquote}
\explab{Explicação}
Aplica-se diretamente à complexidade (entropia) por classe no TCC: classes com mapas mais espalhados não têm necessariamente explicações piores; uma alteração difusa (como a do segmento ST em várias derivações) pode exigir mais regiões que uma localizada.
\end{citacard}

\begin{citacard}{\pagechip{p. 10} §4.2.5 Tendências agregadas}
\begin{origquote}“This behavior suggests that practitioners should not rely on disparity in one metric to serve as a proxy for disparity along all metrics.”\end{origquote}
\begin{tradquote}“Esse comportamento sugere que os profissionais não devem tomar a disparidade em uma métrica como substituta da disparidade em todas as métricas.”\end{tradquote}
\explab{Explicação}
Cada métrica precisa ser comparada entre classes separadamente; uma classe mal posicionada na fidelidade pode estar bem na robustez. É o mesmo cuidado que o TCC aplica ao ler as métricas uma a uma diante do F1.
\end{citacard}

\begin{citacard}{\pagechip{p. 10} §5 Discussão}
\begin{origquote}“As discussed in Section 3, disparity in fidelity and consistency always indicates that the explanation method is amplifying or generating problems. Disparity in stability and sparsity, on the other hand, may either indicate an issue with the explanation method, or an issue with the underlying model to be explained.”\end{origquote}
\begin{tradquote}“Como discutido na Seção 3, a disparidade de fidelidade e de consistência sempre indica que o método de explicação está ampliando ou gerando problemas. Já a disparidade de estabilidade e de esparsidade pode indicar um problema do método de explicação ou um problema do próprio modelo explicado.”\end{tradquote}
\explab{Explicação}
Dá um critério para interpretar as diferenças entre classes no TCC: diferenças na sensibilidade (estabilidade) e na complexidade podem vir do modelo, e por isso faz sentido compará-las com o F1; diferenças na fidelidade e na consistência apontam mais para o método.
\end{citacard}

\begin{citacard}{\pagechip{p. 6} §4.2.1 Disparidade de fidelidade}
\begin{origquote}“These larger variances themselves suggest that prediction gap fidelity is a less robust metric on complex models.”\end{origquote}
\begin{tradquote}“Essas variâncias maiores sugerem, por si, que a \textit{prediction gap fidelity} é uma métrica menos robusta em modelos complexos.”\end{tradquote}
\explab{Explicação}
Em modelos profundos, a variância das métricas dentro de cada grupo pode encobrir diferenças entre grupos. No TCC, com poucas amostras em algumas classes, diferenças de posição entre classes devem ser lidas junto com a dispersão de cada uma.
\end{citacard}

\section{Relação com o TCC}
Entre os trabalhos desta seção, este é o de desenho mais próximo ao do TCC. As quatro propriedades têm correspondentes diretos: fidelidade por perturbação (Faithfulness Correlation e Faithfulness Estimate), estabilidade nos moldes de Yeh et al.\ (Max/Avg-Sensitivity), esparsidade (Complexity) e consistência entre sementes, que é exatamente o \textbf{piso} do TCC (distância entre dois mapas GradientSHAP do mesmo ECG com sementes diferentes). O artigo calcula as métricas por instância e as compara entre grupos, como o TCC faz entre classes, e conclui que as diferenças são frequentes, maiores em modelos complexos e particularmente comuns com Integrated Gradients e SHAP, métodos aparentados ao GradientSHAP. Sua leitura de cada tipo de disparidade (de origem do método ou do modelo) ajuda a interpretar a comparação descritiva com o F1 por classe.

As diferenças: os grupos do artigo são demográficos e binários, e os dados são tabulares; no TCC são nove classes que também são o rótulo do modelo. O artigo usa um teste de Mann-Whitney sobre os valores por instância entre dois grupos, o que é diferente de correlacionar médias de métricas com o F1; um teste desse tipo entre pares de classes seria uma extensão possível, com o cuidado das comparações múltiplas.

\section{Como citar no TCC}
\begin{itemize}
  \item Dai et al.\ (2022) mostraram que fidelidade, estabilidade, consistência e esparsidade das explicações podem diferir de forma significativa entre grupos de dados, sobretudo em modelos não lineares e com métodos como Integrated Gradients e SHAP, e que a disparidade em uma métrica não indica disparidade nas demais; por isso, cada métrica foi comparada separadamente entre as classes.
\end{itemize}

\end{document}
